\documentclass[10pt,a4paper]{article}
\usepackage[utf8]{inputenc}
\usepackage[T1]{fontenc}
\usepackage[a4paper,top=20mm,bottom=20mm,left=16mm,right=16mm,headheight=14pt]{geometry}
\usepackage{fancyhdr}
\usepackage{booktabs}
\usepackage{longtable}
\usepackage{array}
\usepackage{enumitem}
\usepackage{amssymb}
\usepackage{xcolor}
\usepackage[hyphens]{url}
\usepackage[colorlinks=true,urlcolor=blue!60!black,linkcolor=blue!60!black]{hyperref}

\setlength{\parindent}{0pt}
\setlength{\parskip}{3pt}
\renewcommand{\arraystretch}{1.28}

\pagestyle{fancy}
\fancyhf{}
\fancyhead[L]{\small LANGUAGE \& SOCIETY}
\fancyhead[R]{\small TRANSCRIPT ANALYSIS WORKSHEET}
\fancyfoot[C]{\small\thepage}
\renewcommand{\headrulewidth}{0.4pt}

\begin{document}

\begin{center}
    {\large\bfseries LANGUAGE \& SOCIETY PROJECT}\\[3pt]
    {\LARGE\bfseries Transcript Analysis Worksheet}\\[3pt]
    {\normalsize Qualitative analysis of political news discourse}
\end{center}
\vspace{2pt}
\noindent\rule{\textwidth}{0.4pt}
\vspace{4pt}


\section*{Creator 2: Dhruv Rathee (Independent Investigative / Political Explainer YouTuber)}
\begin{itemize}[leftmargin=1.5em, itemsep=3pt, topsep=4pt]
  \item \textbf{Analyst:} Baani Solanki (\texttt{2025114016})
  \item \textbf{Channel / Outlet:} Dhruv Rathee (\texttt{@dhruvrathee}, YouTube)
  \item \textbf{Ideological / Editorial Positioning:} Anti-establishment, progressive-liberal investigative explainer; frames both controversies as systemic regime-level corruption and authoritarian state repression (\texttt{"Tana Shah government"}, \texttt{"dhritarashtra"}) directly implicating Prime Minister Narendra Modi, Home Minister Amit Shah, and the BJP-RSS ecosystem.
  \item \textbf{Transcript Files Referenced:}
  \begin{itemize}[leftmargin=1.4em, itemsep=2pt, topsep=2pt]
    \item Event 1: \texttt{transcripts/event1\_ram\_mandir/02\_dhruv\_rathee\_oeH0fOSJuX8.txt}
    \item Event 2: \texttt{transcripts/event2\_neet\_protest/02\_dhruv\_rathee\_l-6Pm-112io.txt}
    \item \textit{(Caption-track note: Dhruv Rathee's team manually uploaded English-translated Roman-script subtitles preserving key Hindi lexical items such as \texttt{Chanda Chori}, \texttt{Ramlal}, \texttt{dhritarashtra}, \texttt{Anshan}, and \texttt{Tana Shah} directly into YouTube's manual \texttt{hi} native slot; non-AI-dubbed original audio verified.)}
  \end{itemize}
\end{itemize}
\vspace{4pt}\noindent\rule{\textwidth}{0.4pt}\vspace{4pt}

\section*{Guidelines: Five features to look for in every event}
\begin{longtable}{@{}>{\raggedright\arraybackslash}p{0.17\textwidth} >{\raggedright\arraybackslash}p{0.20\textwidth} >{\raggedright\arraybackslash}p{0.24\textwidth} >{\raggedright\arraybackslash}p{0.32\textwidth}@{}}
\toprule
\textbf{Feature} & \textbf{What to look for} & \textbf{Illustrative example} & \textbf{Questions to answer} \\
\midrule
\endhead
\textbf{1. Naming} & What are people, groups, actions, or events called? & "Students", "protesters", "agitators"; "irregularity", "scam", "controversy". & What exact label is used? Who or what does it refer to? Who uses it? What alternative label could have been used? \\
\addlinespace[4pt]
\textbf{2. Presupposition} & What is treated as already true or taken for granted? & "Why are students creating chaos again?" may presuppose earlier disruption. & What does the wording assume? Is that assumption supported, disputed, or left unexamined? Does it appear in a question or words such as "again", "still", or "finally"? \\
\addlinespace[4pt]
\textbf{3. Agency} & Who is shown doing something? Who is acted upon? Who disappears? & "Police injured protesters" vs. "Protesters were injured." & Who is the actor? Who experiences the action? Is responsibility named, backgrounded, or omitted? \\
\addlinespace[4pt]
\textbf{4. Antagonisation} & Who is positioned against whom? & "Ordinary citizens" vs. "political troublemakers". & Which groups are contrasted? Is one presented as more legitimate, moral, patriotic, dangerous, or disloyal? Is an "us vs. them" division implied? \\
\addlinespace[4pt]
\textbf{5. Intensification / softening} & What is exaggerated, minimised, softened, or euphemised? & "Violent mob" vs. "peaceful demonstrators"; "incident" instead of a direct description. & Which words heighten emotion or seriousness? Which minimise it? Are metaphors, dramatic adjectives, hedges, or euphemisms used? Does the wording match the evidence? \\
\bottomrule
\end{longtable}
\textit{Reminder: A label or implication is not automatically false or manipulative. Examine the surrounding context and evidence.}

\vspace{4pt}\noindent\rule{\textwidth}{0.4pt}\vspace{4pt}

\section*{Event1. Ram Mandir donation controversy}
\textbf{Focus:} How is the controversy framed---as financial accountability, religious sentiment, political motivation, or something else? Be precise about the specific allegation or episode covered.

\begin{itemize}[leftmargin=1.5em, itemsep=3pt, topsep=4pt]
  \item \textbf{Outlet:} Dhruv Rathee (YouTube)
  \item \textbf{Date:} \texttt{2026-07-04}
  \item \textbf{Video title:} \textit{Ram Mandir Scam is WORSE Than You Think | Dhruv Rathee} (\url{https://youtu.be/oeH0fOSJuX8?si=tGAGQeFY_MRU32Az})
  \item \textbf{Timestamp(s):} Relevant Event 1 segments \texttt{[00:00--01:54]} and \texttt{[02:49--15:58]} (total video duration \texttt{16:25}; self-promotional AI Masterclass / Bootcamp segments at \texttt{[01:55--02:48]} and \texttt{[15:59--16:25]} excluded per project guidelines).
  \item \textbf{Analyst:} Baani Solanki (\texttt{2025114016})
  \item \textbf{Language:} Hindi / Hinglish (Spoken native Hindi; creator-uploaded Roman-script manual \texttt{hi} track, non-AI-dubbed)
\end{itemize}

\subsection*{Event-specific questions (Event 1)}
\begin{enumerate}[leftmargin=1.6em, itemsep=4pt, topsep=4pt]
  \item \textbf{What terms are used: "scam", "alleged irregularity", "corruption", "political controversy", or something else?}
  \begin{itemize}[leftmargin=1.4em, itemsep=2pt, topsep=2pt]
    \item Rathee uses unmitigated, high-severity crime labels throughout: \texttt{"Scam"} (in title and at \texttt{[01:28]}, \texttt{[08:40]}, \texttt{[11:07]}, \texttt{[13:14]}), \texttt{"Chanda Chori"} / \texttt{"Chanda robbery"} (\texttt{[00:00]}, \texttt{[05:55]}, \texttt{[09:36]}, \texttt{[10:32]}), \texttt{"open-air robbery in the Ram Temple for years"} (\texttt{[09:36]}), \texttt{"dacoits"} (\texttt{[00:00]}, \texttt{[10:32]}), \texttt{"big frauds"} (\texttt{[10:21]}, \texttt{[11:48]}), and \texttt{"looted"} (\texttt{[01:44]}, \texttt{[15:06]}).
  \end{itemize}
  \item \textbf{How are the temple, trust, donors, and people making allegations described?}
  \begin{itemize}[leftmargin=1.4em, itemsep=2pt, topsep=2pt]
    \item \textbf{Trust \& Leadership:} Explicitly stripped of institutional autonomy---\texttt{"this is not an independent Trust. The government is directly involved in this"} (\texttt{[04:46]})---noting that 9 of 15 members have direct Modi/BJP/RSS links, including PM's ex-Principal Secretary Nripendra Misra, CEC Gyanesh Kumar, Anil Mishra, and VHP Vice-President Champat Rai (\texttt{[04:18--05:02]}). Trustees are called \texttt{"temple-owners"} who blocked SBI's warnings (\texttt{[07:33]}), while arrested counting agents (\texttt{Tinnu Yadav}, \texttt{Subhash Srivastav}, \texttt{Lavkush Mishra}, \texttt{K.D. Tiwari}, \texttt{Anukalp Mishra}) are portrayed as low-salaried (\texttt{Rs.\,14,000--20,000/month}) insiders who \texttt{"became billionaires"} (\texttt{Rs.\,50 crore} fortune attributed to Tinnu Yadav at \texttt{[07:33--07:54]}).
    \item \textbf{Donors:} Depicted empathetically as sincere believers whose \texttt{"faith was played with in the name of Lord Ram"} (\texttt{[01:44]})---exemplified by Pune's Sindhi community donating \texttt{210} 1-kg silver bricks worth \texttt{Rs.\,2 crore} (\texttt{[00:17--01:20]})---except for \texttt{"blind devotees who look at Narendra Modi above Lord Ram"} (\texttt{[14:24]}).
    \item \textbf{People Making Allegations:} Whistleblower account-in-charge Mahipal Singh (\texttt{[06:55]}, \texttt{[08:03--08:31]}), SP chief Akhilesh Yadav (\texttt{[08:50]}), Shiv Sena's Sanjay Raut (\texttt{[09:18]}), Mumbai businessman Anil Vishwakarma (\texttt{[09:27]}), BJP politician Rajni Singh (\texttt{[13:14--13:24]}), and AAP's Arvind Kejriwal (\texttt{[13:42]}) are all treated as credible corroborating witnesses.
  \end{itemize}
  \item \textbf{Is the issue framed as financial accountability, an attack on Hindu faith, a political dispute, or a matter of public trust?}
  \begin{itemize}[leftmargin=1.4em, itemsep=2pt, topsep=2pt]
    \item Framed as \textbf{systemic, state-protected financial corruption and betrayal of public/devotee trust} spanning every stage---\texttt{"From the purchase of land, to the construction and donation of the Ram Mandir, there have been frauds"} (\texttt{[11:56]})---enabled by deliberate executive blindness (\texttt{[12:03]}, \texttt{[13:54]}).
  \end{itemize}
  \item \textbf{Are questions about donations presented as legitimate scrutiny or as an attempt to malign the temple or its supporters?}
  \begin{itemize}[leftmargin=1.4em, itemsep=2pt, topsep=2pt]
    \item Presented as \textbf{vital public-interest scrutiny}. When UP CM Yogi Adityanath warns critics to \texttt{"Stop playing with their faith"} (\texttt{[14:13]}), Rathee immediately counters with a rhetorical question: \texttt{"Who is playing with their faith? The scammer or the one who is questioning the scam?"} (\texttt{[14:24]}).
  \end{itemize}
  \item \textbf{Are allegations attributed to specific people and supported with evidence, or presented as established facts?}
  \begin{itemize}[leftmargin=1.4em, itemsep=2pt, topsep=2pt]
    \item Rathee weaves \textbf{named documentary/witness attributions} into an \textbf{established-fact narrative}: he cites whistleblower Mahipal Singh's testimony (\texttt{[06:55]}, \texttt{[08:03--08:40]}), SBI's warning 3 months prior (\texttt{[07:15]}), a November 2020 private audit report (\texttt{[10:12--10:21]}), \textit{Dainik Bhaskar} investigations into \texttt{40\%} construction commissions (Somesh Anand) and \texttt{17x} land-price inflation (\texttt{[10:49--11:24]}), June 2021 land records (\texttt{[11:38]}), and BJP leader Rajni Singh's PMO complaint (\texttt{[13:14--13:42]}). Yet unverified upper-bound figures (\texttt{"donations ranging from Rs 200 crore"}, \texttt{[01:28]}; Tinnu Yadav's \texttt{"fortune of around Rs.\,50 crore"}, \texttt{[07:54]}; Kejriwal's claim of \texttt{"12 IB reports"}, \texttt{[13:42]}) are absorbed smoothly into his indictment.
  \end{itemize}
  \item \textbf{How are denials and counterarguments represented: explained, dismissed, challenged, or omitted?}
  \begin{itemize}[leftmargin=1.4em, itemsep=2pt, topsep=2pt]
    \item \textbf{Quoted briefly only to be dismantled.} Champat Rai's post-June-7 denial (\texttt{"there was no problem in the donation audit"}, \texttt{[08:58]}) is refuted with missing gold/silver receipts (\texttt{[09:08--10:03]}); Champat Rai and Anil Mishra's resignations on 27 June are dismissed as insufficient (\texttt{"Is it enough to resign? If there's a scam here, should they also be arrested?"}, \texttt{[13:05--13:14]}); Yogi Adityanath's counter-attack is inverted (\texttt{[14:13--14:24]}).
  \end{itemize}
  \item \textbf{Does the presenter distinguish criticism of the trust or its financial practices from criticism of the temple or religion itself?}
  \begin{itemize}[leftmargin=1.4em, itemsep=2pt, topsep=2pt]
    \item \textbf{Yes, philosophically, though with a radical anti-donation conclusion.} Rathee invokes the \textit{Swades} poem (\texttt{"Ram is in unity, in progress... Whoever takes out Ravan from the heart, Ram is in his heart"}, \texttt{[15:23--15:40]}) to separate ethical devotion to Lord Ram from \texttt{"the one who shouts Jai Shri Ram the loudest"} while harbouring \texttt{"Ravan"} (\texttt{[15:49]}). However, unlike all other 5 creators, he advises viewers: \texttt{"Whenever you go to a temple, any temple, never donate money... give it to a poor person. Give it to an NGO"} (\texttt{[14:54--15:15]}).
  \end{itemize}
  \item \textbf{Who is portrayed as protecting donors' interests, and who as exploiting religious sentiment?}
  \begin{itemize}[leftmargin=1.4em, itemsep=2pt, topsep=2pt]
    \item \textbf{Exploiting religious sentiment:} Trust General Secretary Champat Rai, his associates/nephews (Tinnu Yadav, Chandan Rai, Somesh Anand), UP CM Yogi Adityanath, and PM Modi/PMO (\texttt{"When you ask for votes in the name of Lord Ram, they stand in the front line first. But when the Lord Ram is being robbed, they let it happen"}, \texttt{[12:03--12:11]}).
    \item \textbf{Protecting donors' interests:} Fired whistleblower Mahipal Singh (\texttt{[08:03--08:31]}), investigative journalists (\textit{Dainik Bhaskar}), dissenting saints like Mahant Kamal Nayan Das (\texttt{[12:30--12:38]}), and opposition/citizen whistleblowers raising pressure (\texttt{[08:50]}, \texttt{[12:11]}, \texttt{[13:42]}).
  \end{itemize}
\end{enumerate}

\subsection*{Apply the five shared features (Event 1)}
\begin{longtable}{@{}>{\raggedright\arraybackslash}p{0.19\textwidth} >{\raggedright\arraybackslash}p{0.38\textwidth} >{\raggedright\arraybackslash}p{0.37\textwidth}@{}}
\toprule
\textbf{Feature} & \textbf{Exact wording + timestamp} & \textbf{What it does in context} \\
\midrule
\endhead
\textbf{1. Naming} & \texttt{"this is not an independent Trust. The government is directly involved... The bank was saying to remove them, but the temple-owners were saying, let them be."} \texttt{[04:46]}, \texttt{[07:33]} \newline  \texttt{"Kejriwal called Modi a dhritarashtra."} \texttt{[13:54]} & Reclassifies fiduciary trustees into proprietary \texttt{"temple-owners"} and fuses the Trust with the Modi government; amplifies the Mahabharata metaphor \texttt{"dhritarashtra"} (the willfully blind king protecting corrupt kin) to frame executive silence as complicity. \\
\addlinespace[4pt]
\textbf{2. Presupposition} & \texttt{"Ram Mandir Scam is WORSE Than You Think"} \texttt{[Title]} \newline  \texttt{"Let's understand how in today's video, people's faith was played with in the name of Lord Ram, and the public was looted."} \texttt{[01:44]} & Both the comparative adjective (\texttt{WORSE Than You Think}) and the \textit{how}-complement (\texttt{understand how... people's faith was played with}) presuppose the existence of a massive scam as indisputable common ground before the main explanation even begins. \\
\addlinespace[4pt]
\textbf{3. Agency} & \texttt{"On the contrary, Mahipal Singh was removed from his job. Apart from that, CCTV footage of 7-8 months was also deleted."} \texttt{[08:22]} vs \texttt{"And the government was blind to this. Was it intentional or was it just like that?"} \texttt{[12:03]} & Uses strategic passives (\texttt{was removed}, \texttt{was deleted}) where insiders covered tracks, then immediately pivots to pointed rhetorical questions forcing the viewer to assign intentional agency up the chain to Champat Rai and the government (\texttt{[08:03]}, \texttt{[12:03]}). \\
\addlinespace[4pt]
\textbf{4. Antagonisation} & \texttt{"the one who shouts Jai Shri Ram the loudest, is not Ram in their hearts, but Ravan. The biggest criminals who rob in the name of Ram, and even more, those who are saving them."} \texttt{[15:49--15:58]} & Constructs a stark mythological-moral binary (\texttt{Ram} vs \texttt{Ravan}) casting vocal Hindutva politicians and Trust protectors not merely as negligent administrators, but as \texttt{"Ravan"} and \texttt{"the biggest criminals"}. \\
\addlinespace[4pt]
\textbf{5. Intensification / softening} & \textbf{Intensification:} \texttt{"Six months passed. Then a year passed. Two years. Five years... everything was disappeared"} \texttt{[01:03--01:28]}; \texttt{"open-air robbery in the Ram Temple for years"} \texttt{[09:36]}; \texttt{"have you ever heard of them happening in one project?"} \texttt{[11:56]} & Narrative pacing (counting elapsed years in short staccato sentences), hyperbole (\texttt{everything was disappeared}, \texttt{became billionaires}), and superlative rhetorical questions intensify moral outrage across the entire 6-year timeline. \\
\bottomrule
\end{longtable}

\subsection*{Short summary (Event 1)}
\begin{itemize}[leftmargin=1.5em, itemsep=3pt, topsep=4pt]
  \item \textbf{What interpretation of the event does this video encourage?} That the June 2026 cash-counting arrests are merely the tip of a six-year, multi-layered institutional plunder---encompassing \texttt{210 kg} missing silver bricks (\texttt{[00:26]}), unrecorded gold/silver jewellery (\texttt{[09:08--10:03]}), \texttt{40\%} construction kickbacks (\texttt{[10:49]}), and \texttt{17x} land-purchase inflation (\texttt{[11:15]})---shielded by BJP-RSS appointees and a willfully blind PMO (\texttt{[04:18--05:22]}, \texttt{[13:14--14:04]}).
  \item \textbf{Whose perspective is foregrounded, challenged, or missing?} \textit{Foregrounded:} Unreceipted donors (Pune Sindhi community), fired whistleblower Mahipal Singh, \textit{Dainik Bhaskar} investigative reporters, and opposition critics. \textit{Challenged:} Champat Rai's denials, the UP SIT's restriction of the FIR to 8 junior staff (\texttt{[12:48--13:05]}), and CM Yogi Adityanath's faith-based deflection (\texttt{[14:13--14:24]}). \textit{Missing:} Any direct interview or detailed rebuttal from Trust members, SBI officials, or SIT investigators explaining why senior trustees were not named in the initial FIR.
  \item \textbf{What evidence supports your interpretation? What complicates it?} \textit{Supports:} Specific donor testimonies (\texttt{210} silver bricks in Jan 2021 at \texttt{[00:17--01:20]}), November 2020 private audit warnings (\texttt{[10:12--10:21]}), Mahipal Singh's on-record account of bundle manipulation (\texttt{12--13} notes hidden in 10-note bundles at \texttt{[06:47--07:06]}) and deleted CCTV footage (\texttt{[08:22--08:40]}), and documented 2021 land-price jumps (\texttt{[11:15--11:38]}). \textit{Complicates:} Several headline numbers (\texttt{Rs.\,200 crore} total disappearance at \texttt{[01:28]}, \texttt{Rs.\,50 crore} for Tinnu Yadav at \texttt{[07:54]}, \texttt{12 IB reports} claimed by Kejriwal at \texttt{[13:42]}) rely on unverified media/political claims, and his sweeping conclusion (\texttt{"never donate money"} to \textit{any} temple, \texttt{[14:54]}) generalizes one Trust's failure to all religious institutions.
\end{itemize}
\vspace{4pt}\noindent\rule{\textwidth}{0.4pt}\vspace{4pt}

\section*{Event2. Students' protest regarding NEET paper leak}
\textbf{Focus:} How does coverage construct the legitimacy of student protest and responsibility for the underlying problem?

\begin{itemize}[leftmargin=1.5em, itemsep=3pt, topsep=4pt]
  \item \textbf{Outlet:} Dhruv Rathee (YouTube)
  \item \textbf{Date:} \texttt{2026-07-24} \textit{(published one day BEFORE Dharmendra Pradhan's resignation on 25 July)}
  \item \textbf{Video title:} \textit{How CJP defeated BJP at Jantar Mantar Protests? | Full Story of 20th July | Dhruv Rathee} (\url{https://youtu.be/l-6Pm-112io?si=TXnmzWqiJKF-qXev})
  \item \textbf{Timestamp(s):} Relevant Event 2 segments \texttt{[00:00--02:39]} and \texttt{[03:01--17:19]} (total video duration \texttt{17:50}; AI Masterclass promo \texttt{[02:40--03:00]} and end-screen promo \texttt{[17:20--17:50]} excluded per project guidelines).
  \item \textbf{Analyst:} Baani Solanki (\texttt{2025114016})
  \item \textbf{Language:} Hindi / Hinglish (Spoken native Hindi; creator-uploaded Roman-script manual \texttt{hi} track, non-AI-dubbed)
\end{itemize}

\subsection*{Event-specific questions (Event 2)}
\begin{enumerate}[leftmargin=1.6em, itemsep=4pt, topsep=4pt]
  \item \textbf{Are participants called "students", "candidates", "protesters", "agitators", "mobs", or "political pawns"?}
  \begin{itemize}[leftmargin=1.4em, itemsep=2pt, topsep=2pt]
    \item Rathee consistently names them \texttt{"protesters"} (\texttt{[00:23]}, \texttt{[01:21]}, \texttt{[08:20]}), \texttt{"the smart people"} (\texttt{[01:21]}), \texttt{"the youth of the country, the people of the country"} (\texttt{[02:55--03:05]}), \texttt{"thousands of students"} (\texttt{[06:33]}, \texttt{[15:41]}), \texttt{"normal people"} (\texttt{[09:26]}), and \texttt{"children"} / \texttt{"its own children"} (\texttt{[16:21--16:44]}), elevating their ethos to \texttt{"the revolutionary spirit, which was seen in the freedom fighters during the independence"} (\texttt{[15:06--15:14]}).
  \end{itemize}
  \item \textbf{Are their stated demands explained, summarised, mocked, questioned, or left out?}
  \begin{itemize}[leftmargin=1.4em, itemsep=2pt, topsep=2pt]
    \item \textbf{Summarised empathetically as minimal and self-evidently just:} \texttt{"the only thing they demanded was the resignation of Education Minister Dharmendra Pradhan"} (\texttt{[03:13]}) and \texttt{"There was only one demand from the children. That the exams should not be stolen. That a better Education Minister should be appointed"} (\texttt{[16:29--16:37]}).
  \end{itemize}
  \item \textbf{Does the presenter focus on the original grievance, disruption, political affiliations, or public inconvenience?}
  \begin{itemize}[leftmargin=1.4em, itemsep=2pt, topsep=2pt]
    \item Focuses overwhelmingly on a \textbf{minute-by-minute forensic reconstruction of pre-meditated state repression on 17--20 July} (\texttt{[00:23--02:19]}, \texttt{[04:27--14:54]}) and nationwide student solidarity (\texttt{[15:14--16:10]}), completely rejecting disruption/inconvenience frames.
  \end{itemize}
  \item \textbf{Who is represented as responsible: examination authorities, government, students, political parties, or others?}
  \begin{itemize}[leftmargin=1.4em, itemsep=2pt, topsep=2pt]
    \item Responsibility is placed squarely on \textbf{PM Narendra Modi} (\texttt{"scared Tana Shah"}, \texttt{[10:28]}; \texttt{"Shame on you, Narendra Modi"}, \texttt{[16:44]}), \textbf{Home Minister Amit Shah} (for abruptly replacing Delhi Police Commissioner Satish Golcha with IB Special Director Anurag Kumar on 17 July at \texttt{[04:27--04:54]}), \textbf{Health Minister J.P. Nadda} (accused of staging a 2-hour distraction meeting to decapitate protest leadership at \texttt{[14:20--14:47]}), and \textbf{Delhi Police / RAF / civil-clothes goons} (\texttt{[05:05]}, \texttt{[11:37--13:25]}).
  \end{itemize}
  \item \textbf{Are questions framed around whether students have legitimate grievances, or why they are causing disruption?}
  \begin{itemize}[leftmargin=1.4em, itemsep=2pt, topsep=2pt]
    \item Framed entirely around \textbf{state conspiracy and abuse of peaceful citizens}: why broken vehicles and brick-laden trucks appeared inside police cordons at \texttt{04:40 AM} (\texttt{[00:23--01:29]}), why metal detectors were removed (\texttt{[01:29--02:01]}), why Commissioner Satish Golcha was sacked (\texttt{[04:27--05:37]}), and why police hid their nameplates and faces (\texttt{[11:37--12:25]}).
  \end{itemize}
  \item \textbf{Are claims of violence, disorder, or political manipulation attributed to sources and substantiated?}
  \begin{itemize}[leftmargin=1.4em, itemsep=2pt, topsep=2pt]
    \item Rathee attributes \textbf{all pre-planned disorder and violence to the state}, substantiating his timeline with timestamped ground videos from independent YouTubers (\texttt{Shyam Meera Singh} at \texttt{[09:03--09:15]}, \texttt{[12:37]}; \texttt{Sarthak Goswami} at \texttt{[10:46]}, \texttt{[12:48]}; \texttt{Sandeesh Bhatia} at \texttt{[12:04]}, \texttt{[13:07]}; \texttt{Being Human / Ramit Verma} bleeding from the head at \texttt{[13:16--13:25]}), DCP New Delhi's \texttt{7:00 PM} Section 163 tweet (\texttt{[05:46--05:55]}), and Sonam Wangchuk's \texttt{11:05 AM} hospital letter (\texttt{[07:52--08:10]}). Conversely, claims of pre-planted brick trucks (\texttt{[00:41--01:29]}) and Amit Shah's motive for transferring Golcha (\texttt{"Perhaps he was upset that the protest was going on peacefully"}, \texttt{[04:36--04:45]}) rely on circumstantial inference.
  \end{itemize}
  \item \textbf{How are police and government actions described compared with protesters' actions? Are similar actions evaluated differently?}
  \begin{itemize}[leftmargin=1.4em, itemsep=2pt, topsep=2pt]
    \item Starkly polarized: Protesters prior to 20 July are \texttt{"100\% non-violent and peaceful"} (\texttt{[00:53]}, \texttt{[03:50]}), and breaking metro/street barricades on 20 July is framed as excusable democratic fervour (\texttt{"When the patience was out of hand, people jumped out of the barriers... People were so excited"}, \texttt{[07:17]}). Police actions, by contrast, are described as a \texttt{"war against the people of its own country"} (\texttt{[02:19]}, \texttt{[16:21]}), \texttt{"kidnapped Sonam Vangchuk"} (\texttt{[05:14]}), \texttt{"Jallianwala Bagh scene. General Dyer"} (\texttt{[10:08]}), and \texttt{"behaving like terrorists"} (\texttt{[11:54]}). Notably, Rathee \textbf{completely omits} the injury of 110--118 police personnel reported by Aaj Tak (\texttt{JXLIs0WtJAM}, \texttt{[18:09]}, \texttt{[19:43]}) and acknowledged by Sarthak Goswami (\texttt{EM1GeZufO6E}, \texttt{[04:37--04:53]}).
  \end{itemize}
  \item \textbf{Are students allowed to explain their demands in their own words, or does the presenter mainly speak for them?}
  \begin{itemize}[leftmargin=1.4em, itemsep=2pt, topsep=2pt]
    \item \textbf{Hybrid narrator-led documentary:} Rathee's scripted voiceover drives the chronological investigation, punctuated by short, emotionally intense ground-video clips of protesters, doctors (\texttt{"We are doctors! We came to help you... They beat us up"}, \texttt{[11:26]}), citizens confronting unbadged police (\texttt{"Where is your tag? Where is your nameplate?"}, \texttt{[11:37--12:25]}), and assaulted YouTubers (\texttt{[12:58--13:25]}).
  \end{itemize}
  \item \textbf{Does the outlet's framing differ between coverage of the NEET paper-leak issue before the protest and coverage once students begin protesting?}
  \begin{itemize}[leftmargin=1.4em, itemsep=2pt, topsep=2pt]
    \item Rathee links the protest directly to his long-standing thesis of creeping authoritarianism (\texttt{"I had exposed the government's taunting of Shahi [tanashahi] two years ago. Today, you can see the result of that"}, \texttt{[17:26]}): once students hit the streets, the issue escalates in his framing from educational incompetence to a full-blown colonial freedom struggle (\texttt{"East India Company"} vs \texttt{"West India Company"} at \texttt{[16:10]}).
  \end{itemize}
\end{enumerate}

\subsection*{Apply the five shared features (Event 2)}
\begin{longtable}{@{}>{\raggedright\arraybackslash}p{0.19\textwidth} >{\raggedright\arraybackslash}p{0.38\textwidth} >{\raggedright\arraybackslash}p{0.37\textwidth}@{}}
\toprule
\textbf{Feature} & \textbf{Exact wording + timestamp} & \textbf{What it does in context} \\
\midrule
\endhead
\textbf{1. Naming} & \texttt{"A Tana Shah government was ready to strangle the people"} \texttt{[02:09]}; \texttt{"kidnapped Sonam Vangchuk in an ambulance"} \texttt{[05:14]}; \texttt{"Narendra Modi, like a scared Tana Shah, leaves the Parliament building and runs away"} \texttt{[10:28]}; \texttt{"West India Company"} \texttt{[16:10]} & Replaces constitutional terms ("Union Government", "preventive detention", "PM") with delegitimizing authoritarian/colonial epithets (\texttt{Tana Shah} [dictator], \texttt{kidnapped}, \texttt{West India Company}---alluding to Gujarat-based corporate/political rule). \\
\addlinespace[4pt]
\textbf{2. Presupposition} & \texttt{"How exactly the youth of the country, the people of the country, were abused and how it was planned, To understand this better, we need to rewind"} \texttt{[02:55--03:05]} & Presupposes as proven fact both that the crackdown was a pre-meditated conspiracy (\texttt{how it was planned}) and that the victims represent the entire nation (\texttt{the youth of the country, the people of the country}). \\
\addlinespace[4pt]
\textbf{3. Agency} & \texttt{"The government started a war against the people of its own country."} \texttt{[02:19]} \newline  \texttt{"the policemen who were mercilessly beating the students, they had hidden their faces with a handkerchief."} \texttt{[11:45--11:54]} vs \texttt{"the barricades had begun to break"} \texttt{[07:17]} & Assigns maximal, intentional active agency to the government/police (\texttt{started a war}, \texttt{mercilessly beating}, \texttt{hidden their faces}), while using an agentless inchoative construction (\texttt{the barricades had begun to break}, \texttt{[07:17]}) that softens protester force against police barriers. \\
\addlinespace[4pt]
\textbf{4. Antagonisation} & \texttt{"Assuming that it was the Jallianwala Bagh scene. General Dyer was giving orders to fire... in a time, Indians were against the East India Company. And today, these Indians are against the West India Company."} \texttt{[10:08]}, \texttt{[16:10]} & Casts an elected Indian government and its security forces as alien colonial oppressors (\texttt{General Dyer}, \texttt{East/West India Company}) and masked \texttt{"terrorists"} (\texttt{[11:54]}) waging war on \texttt{"freedom fighters"} (\texttt{[15:14]}) and \texttt{"its own children"} (\texttt{[16:21]}). \\
\addlinespace[4pt]
\textbf{5. Intensification / softening} & \textbf{Intensification:} \texttt{"100\% non-violent"} \texttt{[00:53]}, \texttt{"very sinister plan"} \texttt{[02:01]}, \texttt{"strangle the people"} \texttt{[02:09]}, \texttt{"terrifying video"} \texttt{[10:08]}, \texttt{"The more the Tana Shahs in history thought themselves to be great, the more terrifying their end would be"} \texttt{[17:00]} & Uses absolute quantifiers (\texttt{100\%}), visceral bodily violence metaphors (\texttt{strangle}), and historical warnings of dictators' downfalls (\texttt{terrifying their end}) to maximize moral indignation, while conjectural hedges (\texttt{Perhaps}, \texttt{[04:36]}, \texttt{[05:37]}) briefly soften unproven motive claims regarding Satish Golcha's transfer. \\
\bottomrule
\end{longtable}

\subsection*{Short summary (Event 2)}
\begin{itemize}[leftmargin=1.5em, itemsep=3pt, topsep=4pt]
  \item \textbf{What interpretation of the event does this video encourage?} That a \texttt{"100\% non-violent"} student sit-in (\texttt{[00:53]}) grew so popular during Sonam Wangchuk's hunger strike (\texttt{[03:40--04:18]}) that a \texttt{"Tana Shah"} Modi-Shah regime (\texttt{[02:09]}, \texttt{[04:27]}) replaced the Delhi Police Commissioner, \texttt{"kidnapped"} Wangchuk (\texttt{[05:14]}), planted stone-filled vehicles to stage a pretext (\texttt{[00:23--01:29]}), shut down 5 metro stations and 150 internet sites (\texttt{[06:23--06:53]}), and unleashed masked/unbadged police and civil-dress goons---only for resilient students across India to reclaim Jantar Mantar that same evening (\texttt{[15:23--16:10]}).
  \item \textbf{Whose perspective is foregrounded, challenged, or missing?} \textit{Foregrounded:} Protesting students, CJP organisers, independent YouTubers assaulted on the ground (\texttt{Shyam Meera Singh}, \texttt{Sarthak Goswami}, \texttt{Sandeesh Bhatia}, \texttt{Ramit Verma}, \texttt{[12:48--13:35]}), and assaulted volunteer doctors (\texttt{[11:26]}). \textit{Challenged:} PM Modi's \texttt{"mother of democracy"} rhetoric (juxtaposed ironically at \texttt{[00:00]} and \texttt{[02:27]}), J.P. Nadda's dialogue offer (\texttt{[14:20--14:47]}), and Delhi Police crowd-control claims. \textit{Missing:} Any mention of the \texttt{110--118} injured police officers (including 5 IPS officers with fractures), opposition party turf wars, or Abhijit Dipke's political background.
  \item \textbf{What evidence supports your interpretation? What complicates it?} \textit{Supports:} Timestamped pre-dawn visuals (\texttt{04:40 AM} broken cars/stone trucks inside cordoned zones at \texttt{[00:23--01:29]}), official Section 163 / DMRC 5-station closure / internet shutdown records (\texttt{[05:55--06:53]}), Commissioner Golcha's abrupt pre-retirement replacement on 17 July (\texttt{[04:27--04:54]}), and clear multi-angle videos of unbadged/masked police and civil-clothed men slapping reporters and beating doctors (\texttt{[11:05--13:25]}). \textit{Complicates:} Total omission of police casualties (\texttt{110+} injured) creates a strictly one-sided casualty ledger, and equating tear-gas/riot positions outside Parliament with \texttt{"Jallianwala Bagh"} and \texttt{"General Dyer"} (\texttt{[10:08]}) relies on hyper-charged historical analogy rather than lethal live-fire evidence.
\end{itemize}
\vspace{4pt}\noindent\rule{\textwidth}{0.4pt}\vspace{4pt}

\section*{Common form: Blank excerpt record}
\textit{Copy or print this page whenever you find a passage worth analysing. Keep the exact wording and enough context to interpret it.}


\subsection*{Excerpt Record 1 (Event 1 --- Ram Mandir Donation Controversy)}
\begin{itemize}[leftmargin=1.5em, itemsep=3pt, topsep=4pt]
  \item \textbf{Outlet:} Dhruv Rathee | \textbf{Date:} \texttt{2026-07-04}
  \item \textbf{Video title:} \textit{Ram Mandir Scam is WORSE Than You Think | Dhruv Rathee}
  \item \textbf{Timestamp(s):} \texttt{[04:18--05:02]} | \textbf{Analyst:} Baani Solanki | \textbf{Language:} Hindi / Hinglish (Roman script)
  \item \textbf{Exact transcript excerpt:}
    \begin{quote}\small\ttfamily [04:18] This trust has a total of 15 members, 6 of whom are religious figures. And most of the other members, have a direct connection with Narendra Modi and BJP RSS. [04:28] For example, Narendra Modi's Principal Secretary, Nipendra Mishra, is the head of this temple construction committee. Similarly, the entire structure of this trust, [04:37] was prepared by the current Chief Election Commissioner, Gyanesh Kumar... [04:46] I'm telling you all this so that you understand that this is not an independent Trust The government is directly involved in this And the people associated with the government are running this temple\end{quote}
  \item \textbf{Feature(s):} $\boxtimes$ Naming $\boxtimes$ Presupposition $\boxtimes$ Agency $\square$ Antagonisation $\square$ Intensification / softening
  \item \textbf{Who is speaking? Is this the outlet's wording or a quotation?} Dhruv Rathee's own scripted analytical narration.
  \item \textbf{What does the wording imply, foreground, soften, or leave unstated?} Explicitly collapses the legal boundary between the Shri Ram Janmabhoomi Teerth Kshetra Trust and the Union Government (\texttt{"this is not an independent Trust"}), ensuring that any financial malpractice inside the temple is attributed upward as state/BJP-RSS culpability.
  \item \textbf{What context or evidence supports or complicates this interpretation?} Supported by the statutory formation of the Trust in February 2020 (\texttt{[04:08]}) and the documented career backgrounds of Nripendra Misra, Gyanesh Kumar, and Champat Rai---a point independently corroborated by Nilanjan Mukhopadhyay on \textbf{The Wire} (\texttt{7rUL8fGgRJA}, \texttt{[11:00--12:27]}). Complicated by \textbf{Aaj Tak} anchor Anjana Om Kashyap's distinction (\texttt{B4L1pz0o4dc}, \texttt{[10:57--11:24]}) that Nripendra Misra chaired only the \textit{Construction Committee} rather than day-to-day \textit{donation counting management}.
  \item \textbf{How does another outlet describe the same event or actor?} \textbf{Republic TV} (\texttt{0H-bruoEkaY}, \texttt{[04:08]}) asserts the exact opposite---\texttt{"there is no religious angle, there is no political angle, there is no government angle in this controversy"}---treating the Trust and the BJP government purely as swift law-enforcers (\texttt{[03:42--03:51]}).
  \item \textbf{My interpretation, with reasons:} Rathee uses institutional network mapping (\texttt{[04:18--05:02]}) early in the video so viewers cannot accept the later FIR (\texttt{[12:48]}) as an isolated failure of 8 low-level counting clerks.
\end{itemize}

\subsection*{Excerpt Record 2 (Event 1 --- Ram Mandir Donation Controversy)}
\begin{itemize}[leftmargin=1.5em, itemsep=3pt, topsep=4pt]
  \item \textbf{Outlet:} Dhruv Rathee | \textbf{Date:} \texttt{2026-07-04}
  \item \textbf{Video title:} \textit{Ram Mandir Scam is WORSE Than You Think | Dhruv Rathee}
  \item \textbf{Timestamp(s):} \texttt{[07:15--08:31]} | \textbf{Analyst:} Baani Solanki | \textbf{Language:} Hindi / Hinglish (Roman script)
  \item \textbf{Exact transcript excerpt:}
    \begin{quote}\small\ttfamily [07:15] 3 months before the arrest, the State Bank of India itself had warned that the counting staff had doubts. That they should be removed... [07:33] The bank was saying to remove them, but the temple-owners were saying, let them be. Why? ... [08:03] According to Mahipal Singh he told Champat Rai about the donation theft in 2021... [08:22] But do you know what happened? On the contrary, Mahipal Singh was removed from his job. Apart from that, CCTV footage of 7-8 months was also deleted.\end{quote}
  \item \textbf{Feature(s):} $\boxtimes$ Naming $\boxtimes$ Presupposition $\boxtimes$ Agency $\boxtimes$ Antagonisation $\boxtimes$ Intensification / softening
  \item \textbf{Who is speaking? Is this the outlet's wording or a quotation?} Rathee's narration combined with reported/embedded testimony from former Trust accounts in-charge Mahipal Singh (\texttt{[08:03--08:40]}).
  \item \textbf{What does the wording imply, foreground, soften, or leave unstated?} Naming Trust administrators \texttt{"the temple-owners"} (\texttt{[07:33]}) implies feudal, unaccountable proprietorship over public donations. Juxtaposing \texttt{"the counting staff"} staying employed against whistleblower Mahipal Singh being \texttt{"removed from his job"} and \texttt{"CCTV footage of 7-8 months... deleted"} implies active cover-up rather than mere administrative oversight.
  \item \textbf{What context or evidence supports or complicates this interpretation?} Supported by Mahipal Singh's on-camera interview clip (\texttt{[08:12--08:40]}) and reports of SBI's prior red flags. Complicated by Nripendra Misra's interview on \textbf{Aaj Tak} (\texttt{B4L1pz0o4dc}, \texttt{[05:09--05:19]}), which blamed \textit{SBI itself} for implementing counting supervision \texttt{"in the most non-serious manner"}.
  \item \textbf{How does another outlet describe the same event or actor?} \textbf{Aaj Tak} panelist Prof. Sangeet Ragi (\texttt{B4L1pz0o4dc}, \texttt{[10:36--10:48]}, \texttt{[12:25--12:56]}) shifts blame onto SBI and portrays Champat Rai as a selfless ascetic (\texttt{"do dhoti aur do kurte par baithne vala aadmi"}) whose integrity cannot be questioned.
  \item \textbf{My interpretation, with reasons:} This excerpt is Rathee's strongest evidentiary pivot from "petty theft by clerks" to "institutional complicity," because firing a whistleblower and deleting 7--8 months of CCTV cannot occur without managerial authority.
\end{itemize}

\subsection*{Excerpt Record 3 (Event 1 --- Ram Mandir Donation Controversy)}
\begin{itemize}[leftmargin=1.5em, itemsep=3pt, topsep=4pt]
  \item \textbf{Outlet:} Dhruv Rathee | \textbf{Date:} \texttt{2026-07-04}
  \item \textbf{Video title:} \textit{Ram Mandir Scam is WORSE Than You Think | Dhruv Rathee}
  \item \textbf{Timestamp(s):} \texttt{[14:13--15:58]} | \textbf{Analyst:} Baani Solanki | \textbf{Language:} Hindi / Hinglish (Roman script)
  \item \textbf{Exact transcript excerpt:}
    \begin{quote}\small\ttfamily [14:13] The CM of UP, Adityanath, is threatening the people who are questioning him. Don't take the fire of the Ram devotees. Stop playing with their faith. [14:24] Stop it. Who is playing with their faith? The scammer or the one who is questioning the scam? There are some blind devotees who look at Narendra Modi above Lord Ram... [14:54] Whenever you go to a temple, any temple, never donate money... [15:49] the one who shouts Jai Shri Ram the loudest, is not Ram in their hearts, but Ravan. The biggest criminals who rob in the name of Ram, and even more, those who are saving them.\end{quote}
  \item \textbf{Feature(s):} $\boxtimes$ Naming $\boxtimes$ Presupposition $\boxtimes$ Agency $\boxtimes$ Antagonisation $\boxtimes$ Intensification / softening
  \item \textbf{Who is speaking? Is this the outlet's wording or a quotation?} Paraphrased warning from UP CM Yogi Adityanath (\texttt{[14:13]}) followed by Dhruv Rathee's direct first-person rebuttal and moral peroration.
  \item \textbf{What does the wording imply, foreground, soften, or leave unstated?} Reverses the speech act of Yogi Adityanath's warning by categorizing it as \texttt{"threatening"} (\texttt{[14:13]}), labels partisan apologists \texttt{"blind devotees"} (\texttt{[14:24]}), and culminates in a mythological inversion (\texttt{Ram} vs \texttt{Ravan}, \texttt{[15:49]}) where state protectors of temple thieves are worse than the thieves themselves (\texttt{"and even more, those who are saving them"}).
  \item \textbf{What context or evidence supports or complicates this interpretation?} Supported by Yogi Adityanath's real Ayodhya speech warning against \texttt{"analagadh tippani"} hurting devotees' sentiments (heard verbatim on \textbf{Aaj Tak} \texttt{B4L1pz0o4dc} at \texttt{[03:02--03:13]}). Complicated by Rathee's sweeping imperative at \texttt{[14:54]} (\texttt{"Whenever you go to a temple, any temple, never donate money"}), which hands critics easy ammunition to accuse him of anti-temple bias.
  \item \textbf{How does another outlet describe the same event or actor?} \textbf{Sarthak Goswami} (\texttt{lb7gZnyZH0c}, \texttt{[06:16--06:26]}) similarly contrasts Lord Ram's character with abusive slogan-chanters (\texttt{"Ram ka charitra chori"}), whereas \textbf{Republic TV} (\texttt{0H-bruoEkaY}, \texttt{[01:12]}, \texttt{[03:22]}) hails Yogi Adityanath's exact same intervention as \texttt{"zero tolerance for any form of corruption"}.
  \item \textbf{My interpretation, with reasons:} Rathee reclaims Hindu ethical vocabulary (\textit{Swades} lyric on Ram vs Ravan) to avoid sounding anti-Hindu while delivering his sharpest anti-BJP and anti-institutional-donation critique.
\end{itemize}

\subsection*{Excerpt Record 4 (Event 2 --- NEET Paper Leak Protest)}
\begin{itemize}[leftmargin=1.5em, itemsep=3pt, topsep=4pt]
  \item \textbf{Outlet:} Dhruv Rathee | \textbf{Date:} \texttt{2026-07-24}
  \item \textbf{Video title:} \textit{How CJP defeated BJP at Jantar Mantar Protests? | Full Story of 20th July | Dhruv Rathee}
  \item \textbf{Timestamp(s):} \texttt{[00:23--02:19]} | \textbf{Analyst:} Baani Solanki | \textbf{Language:} Hindi / Hinglish (Roman script)
  \item \textbf{Exact transcript excerpt:}
    \begin{quote}\small\ttfamily [00:23] On 20th July, at 4.40 am, some people protesting at the Jantar Mantar come out... [00:41] right outside the Jantar Mantar, suddenly, some people saw trucks full of stones... [00:53] everything was 100\% non-violent and peaceful. [01:03] But now, near the nearby Parliament, a broken glass car was also seen standing... [01:29] It was planted to defame the protest by the government. That same night, the Delhi Police removed all the metal detectors and security checks from the Jantar Mantar... [02:01] A very sinister plan was being created behind this... [02:09] A Tana Shah government was ready to strangle the people to save its power. [02:19] The government started a war against the people of its own country.\end{quote}
  \item \textbf{Feature(s):} $\boxtimes$ Naming $\boxtimes$ Presupposition $\boxtimes$ Agency $\boxtimes$ Antagonisation $\boxtimes$ Intensification / softening
  \item \textbf{Who is speaking? Is this the outlet's wording or a quotation?} Dhruv Rathee's investigative narration incorporating an on-site protester's written note (\texttt{"It was planted to defame the protest"}, \texttt{[01:29]}) and ground video (\texttt{[01:41]}).
  \item \textbf{What does the wording imply, foreground, soften, or leave unstated?} Pre-emptively inoculates protesters against \textit{any} subsequent accusation of stone-pelting or vandalism on 20 July by framing stones and damaged cars as a state false-flag operation (\texttt{"sinister plan"}, \texttt{"planted to defame"}), escalated via war/dictatorship metaphors (\texttt{"Tana Shah government"}, \texttt{"started a war"}).
  \item \textbf{What context or evidence supports or complicates this interpretation?} Supported by timestamped \texttt{04:40 AM} ground footage inside heavily cordoned zones and protesters' real-time warnings before clashes began. Complicated by \textbf{Aaj Tak}'s interviews with injured Addl. DCPs Anand Mishra and Jagdev (\texttt{JXLIs0WtJAM}, \texttt{[21:12--23:14]}), who suffered fractures from stones thrown from within the crowd at barricades later that day.
  \item \textbf{How does another outlet describe the same event or actor?} \textbf{Republic TV} (\texttt{HTvb6E\_1Qdc}, \texttt{[01:37--01:47]}) attributes 100\% of street violence to \texttt{"rabid, bloodthirsty mobs"} attacking police and RAF, never mentioning pre-dawn stone trucks or removed metal detectors.
  \item \textbf{My interpretation, with reasons:} Opening the video with \texttt{04:40 AM} pre-dawn anomalies establishes a powerful forensic frame: before showing a single clash at \texttt{[07:43]}, Rathee primes the viewer to interpret all disorder as state sabotage.
\end{itemize}

\subsection*{Excerpt Record 5 (Event 2 --- NEET Paper Leak Protest)}
\begin{itemize}[leftmargin=1.5em, itemsep=3pt, topsep=4pt]
  \item \textbf{Outlet:} Dhruv Rathee | \textbf{Date:} \texttt{2026-07-24}
  \item \textbf{Video title:} \textit{How CJP defeated BJP at Jantar Mantar Protests? | Full Story of 20th July | Dhruv Rathee}
  \item \textbf{Timestamp(s):} \texttt{[10:08--10:28]} \& \texttt{[11:37--12:16]} | \textbf{Analyst:} Baani Solanki | \textbf{Language:} Hindi / Hinglish (Roman script)
  \item \textbf{Exact transcript excerpt:}
    \begin{quote}\small\ttfamily [10:08] Where security forces were taking firing positions. Assuming that it was the Jallianwala Bagh scene. General Dyer was giving orders to fire... [10:28] Meanwhile, Narendra Modi, like a scared Tana Shah, leaves the Parliament building and runs away... [11:37] The most disturbing thing was that the police officers... the nameplates of the vehicles were removed from their uniform. Where is your tag? Where is your nameplate? ... [11:54] they had hidden their faces with a handkerchief. They were behaving like terrorists. Apart from this, we also saw that some people came in ordinary clothes, [12:04] were roaming around with the policemen...\end{quote}
  \item \textbf{Feature(s):} $\boxtimes$ Naming $\square$ Presupposition $\boxtimes$ Agency $\boxtimes$ Antagonisation $\boxtimes$ Intensification / softening
  \item \textbf{Who is speaking? Is this the outlet's wording or a quotation?} Dhruv Rathee's voiceover alternating with direct on-the-ground audio of protesters questioning police (\texttt{"Where is your tag? Where is your nameplate?"}, \texttt{[11:37--11:45]}, \texttt{[12:16]}).
  \item \textbf{What does the wording imply, foreground, soften, or leave unstated?} Performs a radical \textbf{reversal of the "terrorist" label}: whereas pro-government channels call student protesters terrorists, Rathee applies \texttt{"behaving like terrorists"} (\texttt{[11:54]}) and \texttt{"General Dyer"} (\texttt{[10:08]}) to masked, unbadged security personnel, while mocking PM Modi's exit from Parliament as cowardly flight (\texttt{"like a scared Tana Shah... runs away"}, \texttt{[10:28]}).
  \item \textbf{What context or evidence supports or complicates this interpretation?} Unbadged policemen and civil-clothed assailants are corroborated across \textbf{Sarthak Goswami} (\texttt{EM1GeZufO6E}, \texttt{[03:57--04:07]}), \textbf{Mohak Mangal} (\texttt{v-RikSqb0Vw}, \texttt{[01:38]}, \texttt{[04:38]}), \textbf{Aaj Tak} (Akhilesh Yadav quote in \texttt{JXLIs0WtJAM}, \texttt{[06:50--07:08]}), and \textbf{The Wire} (\texttt{pYaTVvHeGoM}, \texttt{[07:34--07:42]}). Even the Jallianwala Bagh visual comparison is independently noted by Naseeruddin Shah on \textbf{The Wire} (\texttt{pYaTVvHeGoM}, \texttt{[10:59--11:08]}). However, security forces taking positions at Parliament's iron gates during a breach attempt (\texttt{[09:59--10:17]}) is standard high-security Parliamentary protocol, complicating the \texttt{"General Dyer"} analogy.
  \item \textbf{How does another outlet describe the same event or actor?} \textbf{Aaj Tak} (\texttt{JXLIs0WtJAM}, \texttt{[20:02--20:41]}) quotes DCP Sachin Sharma describing the exact same deployments neutrally as \texttt{"priventiv arenjment"} around \texttt{"vaital instoleshanans"} [vital installations].
  \item \textbf{My interpretation, with reasons:} Rathee combines undeniable procedural illegality (covered vehicle plates, missing name badges) with colonial/terrorist counter-naming to strip the state apparatus of moral authority.
\end{itemize}

\subsection*{Excerpt Record 6 (Event 2 --- NEET Paper Leak Protest)}
\begin{itemize}[leftmargin=1.5em, itemsep=3pt, topsep=4pt]
  \item \textbf{Outlet:} Dhruv Rathee | \textbf{Date:} \texttt{2026-07-24}
  \item \textbf{Video title:} \textit{How CJP defeated BJP at Jantar Mantar Protests? | Full Story of 20th July | Dhruv Rathee}
  \item \textbf{Timestamp(s):} \texttt{[16:10--17:00]} | \textbf{Analyst:} Baani Solanki | \textbf{Language:} Hindi / Hinglish (Roman script)
  \item \textbf{Exact transcript excerpt:}
    \begin{quote}\small\ttfamily [16:10] This will be an inspiration for many generations to come. That in a time, Indians were against the East India Company. And today, these Indians are against the West India Company. [16:21] Today, the world is witnessing how a country called the Internationalist [Mother of Democracy], conspired against its own children, and waged a war on its own people. [16:29] Do you know what was the strangest part? There was only one demand from the children. That the exams should not be stolen... [16:44] I just want to say this to the Prime Minister, Shame on you, Narendra Modi. The government that is beating up the children, it's time for them to leave... [17:00] The more the Tana Shahs in history thought themselves to be great, the more terrifying their end would be.\end{quote}
  \item \textbf{Feature(s):} $\boxtimes$ Naming $\boxtimes$ Presupposition $\boxtimes$ Agency $\boxtimes$ Antagonisation $\boxtimes$ Intensification / softening
  \item \textbf{Who is speaking? Is this the outlet's wording or a quotation?} Dhruv Rathee speaking directly to the camera and addressing Prime Minister Modi in the second person (\texttt{"Shame on you, Narendra Modi"}, \texttt{[16:44]}).
  \item \textbf{What does the wording imply, foreground, soften, or leave unstated?} Contrasts maximal state aggression (\texttt{"conspired against its own children, and waged a war"}) with minimal, innocent youth agency (\texttt{"only one demand from the children. That the exams should not be stolen"}), culminating in a direct political prophecy of regime fall (\texttt{"time for them to leave"}, \texttt{"terrifying their end"}).
  \item \textbf{What context or evidence supports or complicates this interpretation?} Remarkably, \textbf{the very next day (25 July 2026)} Education Minister Dharmendra Pradhan resigned and the government accepted CJP's core demands (confirmed in Sarthak \texttt{EM1GeZufO6E} \texttt{[00:08]} and Republic \texttt{HTvb6E\_1Qdc} \texttt{[00:00]}), showing immense public pressure---yet also complicating Rathee's \texttt{"Tana Shah"} thesis because an unyielding dictatorship rarely concedes ministerial resignation within 5 days of a march.
  \item \textbf{How does another outlet describe the same event or actor?} \textbf{Republic TV} (\texttt{HTvb6E\_1Qdc}, \texttt{[01:56--02:56]}) spins Pradhan's resignation the next day not as a defeat of the BJP, but as a masterstroke whereby \texttt{"the government outplayed"} Rahul Gandhi and ended the street crisis.
  \item \textbf{My interpretation, with reasons:} Rathee transforms a single-issue exam-reform protest into an existential plebiscite on Modi's rule (\texttt{"West India Company"}), mobilizing youth grievance through anti-colonial and familial (\texttt{"children"}) framing.
\end{itemize}
\vspace{4pt}\noindent\rule{\textwidth}{0.4pt}\vspace{4pt}

\section*{Common form: End-of-event summary}

\subsection*{End-of-Event Summary --- Event 1 (Ram Mandir Donation Controversy)}
\begin{itemize}[leftmargin=1.5em, itemsep=3pt, topsep=4pt]
  \item \textbf{The main way this outlet framed the event was:} As a comprehensive, six-year institutional plunder (\texttt{"From the purchase of land, to the construction and donation of the Ram Mandir, there have been frauds"}, \texttt{[11:56]}) run by a government-linked BJP-RSS Trust that ignored 2020 private audit and SBI warnings, fired whistleblower Mahipal Singh, deleted 7--8 months of CCTV, and scapegoated 8 low-level counting staff while shielding senior figures like Champat Rai under a \texttt{"dhritarashtra"} PMO (\texttt{[04:18--08:40]}, \texttt{[12:48--14:04]}).
  \item \textbf{The strongest linguistic evidence was (quote + timestamp):} \texttt{"this is not an independent Trust. The government is directly involved... the one who shouts Jai Shri Ram the loudest, is not Ram in their hearts, but Ravan. The biggest criminals who rob in the name of Ram, and even more, those who are saving them."} (\texttt{[04:46]} \& \texttt{[15:49--15:58]}).
  \item \textbf{The shared feature(s) most clearly illustrated were:} \textbf{Presupposition} (\texttt{Scam is WORSE Than You Think}; \texttt{"understand how... the public was looted"}, \texttt{[01:44]}), \textbf{Naming} (\texttt{temple-owners} \texttt{[07:33]}, \texttt{dhritarashtra} \texttt{[13:54]}, \texttt{Ravan} \texttt{[15:49]}), and \textbf{Antagonisation} (sincere devotees \& whistleblowers vs hypocritical \texttt{"Jai Shri Ram"} scammers and \texttt{"blind devotees"}, \texttt{[14:24--15:58]}).
  \item \textbf{An alternative interpretation of the same evidence is:} That the cash-counting theft (\texttt{\textasciitilde{}Rs.\,80 lakh} verified on CCTV in Republic \texttt{0H-bruoEkaY} \texttt{[07:26]}) resulted from lax operational protocols between SBI and counting contractors (as Nripendra Misra admitted on Aaj Tak \texttt{B4L1pz0o4dc} \texttt{[04:35--06:05]}), and that arresting 8 direct culprits within days alongside Champat Rai and Anil Mishra's resignations (\texttt{[12:56--13:05]}) demonstrated swift corrective accountability rather than PMO-directed collusion.
  \item \textbf{What I cannot conclude from this transcript alone:} Whether the unverified upper-end allegations (\texttt{Rs.\,200 crore} missing donations at \texttt{[01:28]}, \texttt{Rs.\,50 crore} amassed by Tinnu Yadav at \texttt{[07:54]}, or \texttt{12 IB reports} sent to the PM at \texttt{[13:42]}) are backed by forensic bank/SIT records, or whether senior trustees personally received kickbacks.
\end{itemize}

\subsection*{End-of-Event Summary --- Event 2 (Students' Protest Regarding NEET Paper Leak)}
\begin{itemize}[leftmargin=1.5em, itemsep=3pt, topsep=4pt]
  \item \textbf{The main way this outlet framed the event was:} As a pre-meditated authoritarian \texttt{"war against the people of its own country"} (\texttt{[02:19]}) in which a \texttt{"Tana Shah"} government (\texttt{[02:09]}) sacked the Delhi Police Commissioner, \texttt{"kidnapped"} hunger-striker Sonam Wangchuk (\texttt{[05:14]}), planted stone trucks (\texttt{[00:41--01:29]}), imposed an internet/metro siege (\texttt{[06:12--06:53]}), and unleashed masked/unbadged police and civil-dress goons on \texttt{"100\% non-violent"} \texttt{"children"}---who nevertheless triumphed like anti-colonial freedom fighters (\texttt{[15:06--16:37]}).
  \item \textbf{The strongest linguistic evidence was (quote + timestamp):} \texttt{"A Tana Shah government was ready to strangle the people to save its power. The government started a war against the people of its own country... Indians were against the East India Company. And today, these Indians are against the West India Company."} (\texttt{[02:09--02:19]} \& \texttt{[16:10]}).
  \item \textbf{The shared feature(s) most clearly illustrated were:} \textbf{Naming} (\texttt{Tana Shah}, \texttt{kidnapped}, \texttt{General Dyer}, \texttt{behaving like terrorists}, \texttt{West India Company}), \textbf{Agency} (attributing full pre-planned agency to Modi/Shah/Delhi Police while rendering broken barricades agentless at \texttt{[07:17]} and omitting all \texttt{110+} police injuries), and \textbf{Intensification} (\texttt{100\% non-violent}, \texttt{strangle}, \texttt{waged a war on its own children}).
  \item \textbf{An alternative interpretation of the same evidence is:} That closing nearby metro stations and erecting Section 163 barricades around Parliament during a Monsoon Session (\texttt{[05:55--06:53]}) were preventive security measures for a high-security constitutional zone, which escalated into violent two-way clashes once lakhs of protesters broke welded barricades (\texttt{[07:17--08:20]}) and pelted stones injuring over 110 police personnel (Aaj Tak \texttt{JXLIs0WtJAM}, \texttt{[19:43--23:14]}).
  \item \textbf{What I cannot conclude from this transcript alone:} Who actually parked the broken vehicles/brick trucks filmed at \texttt{04:40 AM} (\texttt{[00:23--01:12]}), why Satish Golcha was administratively transferred on 17 July (\texttt{[04:27--04:45]}), or how many police officers and protesters respectively initiated physical force at each flashpoint.
\end{itemize}
\textit{Research reminder: interpret wording in context, distinguish allegations from established facts, and support claims with exact excerpts. Mixed or contradictory evidence is worth recording.}


\end{document}
